% ----------------------------------------------------------------------
\section{Optimisation continue avec contraintes}
% ----------------------------------------------------------------------

% ----------------------------------------------------------------------
\begin{frame}<beamer>
  \frametitle{Sommaire}
  \tableofcontents[currentsection]
\end{frame}

% ----------------------------------------------------------------------
\begin{frame}
  \frametitle{Formulation mathématique}

  \[
  (\text{P}^\bullet) \left\{
  \begin{array}{c}
    \text{trouver} \ x \ \text{qui minimise} \ f(x) \ \text{tel que :} \\
    g_i(x) \leq 0, \ i \in I = \{1, \dots, m\} \\
    x \in \R^n
  \end{array}
  \right.
  \]

  \begin{block}{Le problème peut}
    \begin{itemize}
    \item être \emph{infaisable} :
      $X := \{ x \in \R^n \ | \ g_i(x) \leq 0, i \in I \}$,
      l'ensemble des solutions candidates est vide.
    \item être \emph{non borné} : $\exists t, v, x^0$, tels que $\lim_{t \rightarrow \infty} f(x^0 + tv) = -\infty$.
    \item avoir une ou \emph{plusieurs solutions}. 
    \end{itemize}
  \end{block}
  
\end{frame}

% ----------------------------------------------------------------------
\begin{frame}
  \frametitle{Problème infaisable}

  \begin{center}
      \includegraphics[width=0.5\textwidth]{p-infaisable}    
  \end{center}
  
\end{frame}

% ----------------------------------------------------------------------
\begin{frame}
  \frametitle{Problème non borné}

  \begin{center}
    \includegraphics[width=0.4\textwidth]{p-non-borne}    
  \end{center}

\end{frame}

% ----------------------------------------------------------------------
\begin{frame}
  \frametitle{Plusieurs solutions}

  \begin{center}
    \includegraphics[width=0.4\textwidth]{p-plusieurs}    
  \end{center}
  
\end{frame}


%% % ----------------------------------------------------------------------
%% \begin{frame}
%%   \frametitle{Saturation et cône tangent en un point $x^0 \in X$}

%%   \begin{itemize}
%%     \item Une contrainte $g_i$ est \emph{saturée} en $x^0$ ssi $g_i(x^0) = 0$
%%     \item L'ensemble des contraintes saturées en $x^0$ est
%%       $\{g_i\}_{i \in I_0}, \ I_0 := \{i \ | \ g_i(x^0) = 0 \}$
%%     \item Le cône tangent en $x^0$ est donné par le gradient des contraintes saturées en $x^0$ :
%%       $G_0 := \{ y \in \R^n \ | \ {\nabla g_i}^T \cdot y \leq 0, \ \forall i \in I_0 \}$ 
%%   \end{itemize}

%%   \begin{center}
%%       \includegraphics[width=0.45\textwidth,page=1]{cone-tangent} \hspace{0.01\textwidth}   
%%       \includegraphics[width=0.45\textwidth,page=2]{cone-tangent}    
%%   \end{center}
  
%% \end{frame}

%% % ----------------------------------------------------------------------
%% \begin{frame}
%%   \frametitle{Hypothèse de qualification des contraintes (QC)}

%%   $X$ satisfait en $x^0 \in X$ l'hypothèse (QC) ssi pour tout $\epsilon$
%%   positif suffisamment petit, la boule $B(x^0,\epsilon)$ de centre $x^0$
%%   et rayon $\epsilon$ est telle que $B(x^0,\epsilon) \cap G_0 \subset X$. 

%%   \only<1>{
%%   \begin{exampleblock}{Cas où (QC) est vérifiée en $x^0$ : }
%%   \begin{center}
%%       \includegraphics[width=0.5\textwidth,page=1]{cone-tangent}    
%%   \end{center}
%%   \end{exampleblock}
%%   }
%%   \only<2>{
%%   \begin{exampleblock}{Cas où (QC) n'est pas vérifiée en $x^0$ : }
%%   \begin{center}
%%       \includegraphics[width=0.5\textwidth,page=2]{cone-tangent}    
%%   \end{center}
%%   \end{exampleblock}
%%   }
%%   \only<3>{
%%     \begin{block}{Théorème local}
%%       $(QC)$ est vérifiée \alert{en un point $x^0 \in X$}
%%       \begin{itemize}
%%       \item si les gradients $\{ {\nabla g_i}(x^0) \}_{i \in I_0}$ des contraintes saturées en $x^0$
%%         sont linéairement indépendants
%%       \end{itemize}
%%     \end{block}
    
%%     \begin{block}{Théorème global}
%%       $(QC)$ est vérifiée \alert{pour tout $x \in X$}
%%       \begin{itemize}
%%       \item si toutes les fonctions $g_i$ sont linéaires
%%       \item ou si toutes les fonctions $g_i$ sont convexes, différentiables, et l'intérieur de $X$ n'est pas vide
%%       \end{itemize}
%%     \end{block}
%%   }
  
%% \end{frame}

% ----------------------------------------------------------------------
\begin{frame}
  \frametitle{Intuition sur les conditions d'optimalité}

  On ne peut pas faire mieux sans violer les contraintes

  ~
  
  
  \begin{center}
      \includegraphics[width=0.5\textwidth]{kkt}    
  \end{center}

\end{frame}

% ----------------------------------------------------------------------
\begin{frame}
  \frametitle{Hypothèses techniques}

  \[
  (\text{P}^\bullet) \left\{
  \begin{array}{c}
    \text{trouver} \ x \ \text{qui minimise} \ f(x) \ \text{tel que :} \\
    g_i(x) \leq 0, \ i \in I = \{1, \dots, m\} \\
    x \in \R^n
  \end{array}
  \right.
  \]

  \begin{itemize}
  \item (H1) $f$ et $(g_i)_{i \in I}$ continues, différentiables
  \item (H2) $X \neq \emptyset$
  \item (QC) hypothèse de qualification des contraintes
  \end{itemize}
  
\end{frame}

% ----------------------------------------------------------------------
\begin{frame}
  \frametitle{Conditions d'optimalité de Karush, Kuhn et Tucker}

  \begin{block}{Conditions nécessaires}
    En supposant (H1), (H2), (QC), si $x^0$ est un minimum local de $(\text{P}^\bullet)$,
    alors il existe $\{\lambda_i \geq 0\}_{i \in I}$ tels que:
    \[
    (\text{KKT})
    \left\{
    \begin{array}{ll}
      -{\nabla f}(x^0) = \sum_{i \in I} \lambda_i {\nabla g_i}(x^0)
      & \text{(stationnarité)} \\
      \lambda_i g_i(x^0) = 0, \ \forall i \in I
      & \text{(complémentarité)} \\
    \end{array}
    \right.
    \]
  \end{block}

  \begin{block}{Conditions suffisantes}
    Dans le cas où, en plus, $f$ et $\{g_i\}_{i\in I}$ sont \emph{convexes},
    si les conditions (KKT) sont vérifiées pour $x^0$, alors $x^0$ est un minimum
    (global) de $(\text{P}^\bullet)$.  
  \end{block}
  
\end{frame}

% ----------------------------------------------------------------------
\begin{frame}
  \frametitle{Exemple numérique}

  \begin{center}  
  \begin{tabular}{l}
    trouver $(x_1,x_2)$ qui minimisent \\
    \textcolor{red}{$f(x_1,x_2) := (x_1 - 3)^2 + (x_2 - 4)^2$} tels que : \\
    pour tout $i = 1,2,3$, \textcolor{blue}{$g_i(x_1,x_2) \leq 0$}, avec \\
    $g_1(x_1,x_2) := x_1 + x_2 - 5$, $g_2(x_1,x_2) := -x_1$, $g_3(x_1,x_2) := -x_2$ 
    %% \textcolor{blue}{$\left\{
    %% \begin{array}{l}
    %%   g_1(x_1,x_2) \leq 0, \ g_1(x_1,x_2) := x_1 + x_2 - 5 \\
    %%   g_2(x_1,x_2) \leq 0, \ g_2(x_1,x_2) := x_1 - x_2 - 2.5 \\
    %%   g_3(x_1,x_2) \leq 0, \ g_3(x_1,x_2) := -x_1 \\
    %%   g_4(x_1,x_2) \leq 0, \ g_4(x_1,x_2) := -x_2
    %% \end{array}
    %% \right.$}
  \end{tabular}
  \end{center}
  
  \begin{center}  
    \includegraphics[width=0.45\textwidth,page=1]{kuhn-tucker}
  \end{center}
\end{frame}

% ----------------------------------------------------------------------
\begin{frame}
  \frametitle{Résolution par les équations (KKT)}

  \only<1>{  
  \begin{itemize}
    \item Comme pour tout $x$, $\nabla^2f(x)$ est définie-positive, $f$ est convexe
    \item Les contraintes $\{g_i\}_{i \in I}$ sont linéaires donc convexes
    \item $x^\star$ et $\lambda_1,\lambda_2,\lambda_3$ qui vérifient (KKT)
      sont tels que $x^\star$ est le minimum global
    \end{itemize}
  }
  \only<2>{
    On joue aux devinette et suppose que seule $g_1$ est saturée 
  \begin{itemize}
  \item c-à-d. telle que $g_1(x^\star) = x_1^\star + x_2^\star - 5 = 0$,
  \item complémentarité : $\lambda_2 = \lambda_3 = 0$, 
  \item stationnarité : $-{\nabla f}(x^\star) = \lambda_1 {\nabla g_1}(x^\star)$.
  \end{itemize}

  Le système à résoudre est donc: 
  \[
  \left\{
  \begin{array}{l}
    x_1^\star + x_2^\star - 5 = 0 \\
  - \left( \begin{array}{c} 2x_1^\star - 6 \\ 2x_2^\star - 8 \end{array} \right) =
  \lambda_1 \left( \begin{array}{c} 1 \\ 1 \end{array} \right) \\
  \end{array}
  \right.
  \]
  D'où on déduit $x_1^\star = 2$, $x_2^\star = 3$, $\lambda_1 = 2$. 
  }

  \only<3>{
  \begin{center}  
    \includegraphics[width=0.45\textwidth,page=2]{kuhn-tucker}
  \end{center}
  }
\end{frame}

% ----------------------------------------------------------------------
\begin{frame}
  \frametitle{Méthodes d'optimisation}

  \begin{itemize}
  \item Recherche combinatoire et résolution des équations (KKT)
    %par des méthodes d'optimisation sans contrainte
    %ex. resoudre ax = b revient a minimiser (ax - b)^2  
  \item Méthodes des pénalités
    %: on élimine les contraintes en ajoutant à la fonction objectif des pénalités qui assurent
    %qu'un point $\bar{x} \notin X$ ne soit pas minimum   
  \item Méthodes exploitant la théorie de la dualité de Lagrange
    %algorithmes d'Uzawa, Arrow-Hurwicz, Dantzig, méthode des multiplicateurs 
  \item Méthodes de descente sous contrainte
    %: de la solution courante, on cherche un déplacement de descente vers un point de $X$
    %méthodes des directions réalisables, du gradient projeté, du gradient réduit (généralisé)...
  \item Méthodes de linéarisation
    %: à chaque étape, on cherche le minimum d'une approximation linéaire de $f$ pour trouver le déplacement 
    %méthode de Franck et Wolf, plans sécants de Kelley, génération de colonne de Dantzig...
  \item \dots
  \end{itemize}
  
\end{frame}

% ----------------------------------------------------------------------
\begin{frame}
  \frametitle{Note : cas particulier de fonctions linéaires}

  On appelle problème d'optimisation linéaire (ou \emph{programme linéaire}),
  le problème d'optimisation suivant : 
  \[
  \text{(PL)} \left\{
  \begin{array}{c}
    \text{trouver} \ x \ \text{qui minimise} \ f(x) \ \text{tel que :} \\
    g_i(x) \leq 0, \ i \in I = \{1, \dots, m\} \\
    x \in \R^n \\
    f, \{g_i\}_{i \in I} \ \text{sont linéaires} \\
  \end{array}
  \right.
  \]

  (PL) est résolu en pratique par la méthode du \emph{simplexe}
  %% dû essentiellement à Dantzig (1947). Le traitement correct
  %% des configurations dégénérées est dû à Bland (1977).
  %% cf. https://en.wikipedia.org/wiki/Revised_simplex_method
  ou celle \emph{des points intérieurs}.
\end{frame}
